\documentclass[UTF8]{ctexart}
\usepackage{xeCJK}
\usepackage{geometry}
\geometry{left=2cm,right=2cm,top=2.5cm,bottom=2.5cm}
\setlength{\headheight}{12.7pt}

\usepackage{titlesec}
\titleformat{\section}
  {\normalfont\large\bfseries}
  {\thesection}
  {1em}
  {}
\titlespacing*{\section}
  {0pt}
  {3.5ex plus 1ex minus .2ex}
  {2.3ex plus .2ex}

\usepackage{amsmath}
\usepackage{amssymb}
\usepackage{amsthm}
\usepackage{enumitem}
\setlist[enumerate]{itemsep=0pt,topsep=0pt,parsep=0pt,leftmargin=0pt,itemindent=2.5em}
\setlist[itemize]{itemsep=0pt,topsep=0pt,parsep=0pt,leftmargin=0pt,itemindent=2.5em}
\usepackage{fancyhdr}
\pagestyle{fancy}
\fancyhf{}
\usepackage{graphicx}
\usepackage{hyperref}
\usepackage{indentfirst}
\usepackage{lmodern}


\begin{document}
\setlength{\abovedisplayskip}{1pt}
\setlength{\belowdisplayskip}{1pt}

\section{序数指数}

\hypertarget{sec:定义1.1}{}
\noindent\textbf{定义1.1 (序数指数)}

给定非零\href{../02 序数/01 序数#nameddest=sec:定义1.1}{序数} $\alpha$,
递归定义类映射$\mathbf{Exp}_\alpha:\mathbf{On}\to\mathbf{On}$如下:
\vspace{3pt}
\[
    \forall\beta\in\mathbf{On}:\quad
    \mathbf{Exp}_\alpha(\beta)\overset{\mathrm{def}}{=}\left\{\begin{aligned}
        &1,\quad&&\beta=0,\\[-3pt]
        &\mathbf{Exp}_\alpha(\max\beta)\cdot\alpha,\quad&&
        \beta\text{\,\,是后继序数},\\[-3pt]
        &\sup\mathbf{Exp}_\alpha[\beta],\quad&&
        \beta\text{\,\,是极限序数}.
    \end{aligned}\right.
\]

接下来, 对任意的非零序数$\alpha,\beta$, 定义
\[
    \alpha^\beta\overset{\mathrm{def}}{=}\mathbf{Exp}_\alpha(\beta),
\]
称为\textbf{序数指数}. 另外, 规定$0^0=1$, 对任意的非零序数$\beta$有$0^\beta=0$.

\vspace{10pt}

\hypertarget{sec:定理1.1}{}
\noindent\textbf{定理1.1 (序数指数)}

任给序数$\alpha,\beta$,
\begin{enumerate}
    \item $\alpha^0=1$, $\alpha^1=\alpha$,
    \item $\alpha^{\beta^+}=\alpha^\beta\alpha$,
    \item 如果$\alpha$非零且$\beta$是极限序数,
    那么$\displaystyle\alpha^\beta=\sup_{\gamma<\beta}\alpha^\gamma$.
\end{enumerate}

\vspace{15pt}

显然\href{../../03 自然数/02 自然数算术/03 自然数的指数#nameddest=sec:定义1.1}
{自然数的指数}是序数指数的特例.
易证任给序数$\alpha$有$1^\alpha=1,2^\alpha\geqslant\alpha$.

\vspace{10pt}

\hypertarget{sec:定理1.2}{}
\noindent\textbf{定理1.2}

任给序数$\alpha,\beta,\beta'$, 如果$\alpha>1$且$\beta<\beta'$,
那么$\alpha^\beta<\alpha^{\beta'}$.

\noindent{\kaishu 证明.}

给定$\alpha>1$和$\beta$, 对$\beta'>\beta$归纳证明$\alpha^\beta<\alpha^{\beta'}$.
任取$\beta'>\beta$, 假设
\[
    \forall\gamma:\quad\beta<\gamma<\beta'\to\alpha^\beta<\alpha^\gamma.
\]

\begin{enumerate}
    \item 如果$\beta'$是后继序数, 即有$\beta'=\gamma^+$,
    那么有$\beta\leqslant\gamma<\beta'$, 依假设有
    \[
        \alpha^\beta\leqslant\alpha^\gamma<\alpha^{\gamma^+}=\alpha^{\beta'}.
    \]
    \item 如果$\beta'$是极限序数, 则$\beta^+<\beta'$, 从而
    \begin{equation*}
        \alpha^\beta<\alpha^{\beta^+}\leqslant
        \sup_{\gamma<\beta'}\alpha^\gamma=\alpha^{\beta'}.
    \tag*{\qedsymbol}\end{equation*}
\end{enumerate}

\vspace{10pt}

\hypertarget{sec:定理1.3}{}
\noindent\textbf{定理1.3}

任给序数$\alpha,\alpha',\beta$, 如果$\alpha<\alpha'$,
那么$\alpha^\beta\leqslant(\alpha')^\beta$.

\noindent{\kaishu 证明.}

给定$\alpha<\alpha'$,不妨设$\alpha$非零,
对$\beta$归纳证明$\alpha^\beta\leqslant(\alpha')^\beta$.
任取$\beta$, 假设对任意的$\gamma<\beta$有$\alpha^\gamma\leqslant(\alpha')^\gamma$.

\begin{enumerate}
    \item 如果$\beta=0$, 显然.
    \item 如果$\beta$是后继序数, 即有$\beta=\gamma^+$, 那么$\gamma<\beta$, 依假设有
    \[
        \alpha^\beta=\alpha^\gamma\alpha\leqslant
        (\alpha')^\gamma\alpha'=(\alpha')^\beta.\\[1pt]
    \]
    \item 如果$\beta$是极限序数, 那么依假设有
    \vspace{-1pt}
    \begin{equation*}
        \alpha^\beta=\sup_{\gamma<\beta}\alpha^\gamma\leqslant
        \sup_{\gamma<\beta}(\alpha')^\gamma=(\alpha')^\beta.
    \tag*{\qedsymbol}\end{equation*}
\end{enumerate}





\newpage

\vspace{10pt}

\hypertarget{sec:定理1.4}{}
\noindent\textbf{定理1.4}

任给序数$\alpha,\beta$, 如果$\alpha\neq 0$且$\beta>1$,
那么存在唯一的一组序数$\nu,\eta,\varepsilon$使得
\[
    \alpha=\beta^\nu\eta+\varepsilon\quad\land\quad
    0<\eta<\beta\quad\land\quad\varepsilon<\beta^\nu.
\]

\noindent{\kaishu 证明.}

\noindent{\kaishu 存在性.}

取满足$\beta^\rho>\alpha$的最小的$\rho$,
即对任意的$\nu<\rho$有$\beta^\nu\leqslant\alpha$. 显然$\rho$非零.

\vspace{3pt}
如果$\rho$是极限序数, 那么$\displaystyle
\beta^\rho=\sup_{\nu<\rho}\beta^\nu\leqslant\alpha$矛盾.
所以$\rho$是后继序数, 有$\rho=\nu^+$.

\vspace{3pt}
记$\eta=\alpha/\beta^\nu,\varepsilon=\alpha\%\beta^\nu$, 则有
\[
    \alpha=\beta^\nu\eta+\varepsilon\quad\land\quad\varepsilon<\beta^\nu,
\]
同时
\[
    \beta^\nu\eta\leqslant\beta^\nu\eta+\varepsilon=\alpha<\beta^\rho=
    \beta^\nu\beta\quad\implies\quad\eta<\beta,
\]
最后假设$\eta=0$则$\alpha=\varepsilon<\beta^0=1$矛盾, 所以$\eta>0$.

\noindent{\kaishu 唯一性.}

假设同时存在$\nu,\eta,\varepsilon$和$\nu',\eta',\varepsilon'$使得
\[\left\{\begin{aligned}
    &\alpha=\beta^\nu\eta+\varepsilon\,&&\land\quad
    0<\eta<\beta\,&&\land\quad\varepsilon<\beta^\nu,\\[-3pt]
    &\alpha=\beta^{\nu'}\eta'+\varepsilon'\,&&\land\quad
    0<\eta'<\beta\,&&\land\quad\varepsilon'<\beta^{\nu'}.
\end{aligned}\right.\]
假设$\nu<\nu'$, 即$\nu^+\leqslant\nu'$, 则
\[
    \alpha=\beta^\nu\eta+\varepsilon<\beta^\nu\eta+\beta^\nu=
    \beta^\nu\eta^+\leqslant\beta^\nu\beta=\beta^{\nu^+}\leqslant
    \beta^{\nu'}\leqslant\beta^{\nu'}\eta'+\varepsilon'=\alpha,
\]
矛盾. 同理$\nu'<\nu$也矛盾, 所以$\nu=\nu'$.
进一步显然$\eta=\eta',\,\varepsilon=\varepsilon'$. \hfill $\qedsymbol$

\vspace{10pt}

\hypertarget{sec:定理1.5}{}
\noindent\textbf{定理1.5}

任给序数$\alpha>1,\,\beta\neq 0$ :
\begin{enumerate}
    \item 如果$\alpha$是后继序数且$\beta<\omega$, 那么$\alpha^\beta$是后继序数,
    \item 如果$\alpha$是极限序数或$\beta\geqslant\omega$,
    那么$\alpha^\beta$是极限序数.
\end{enumerate}

\noindent{\kaishu 证明.}

\begin{enumerate}
    \item 对$0<\beta<\omega$归纳证明$\alpha^\beta$与$\alpha$同类.
    
    \hspace{2.17em}
    首先$\alpha^1=\alpha$与$\alpha$同类;
    其次任取自然数$\beta$, 如果$\alpha^\beta$与$\alpha$同类, 那么
    \[
        \alpha^{\beta^+}=\alpha^\beta\alpha
    \]
    也与$\alpha$同类.

    \item 对$\beta\geqslant\omega$归纳证明$\alpha^\beta$是极限序数.
    任取$\beta\geqslant\omega$, 假设
    \[
        \forall\gamma:\quad\omega\leqslant\gamma<\beta\to
        \alpha^\gamma\text{\,是极限序数\,}.
    \]

    \hspace{2.17em}
    如果$\beta$是后继序数, 即有$\beta=\gamma^+$,
    那么显然$\omega\leqslant\gamma<\beta$, 依假设有
    \vspace{3pt}
    \[
        \alpha^\gamma\text{\,是极限序数\,}\to\,
        \alpha^\beta=\alpha^\gamma\alpha\text{\,是极限序数\,}.
    \]

    \vspace{3pt}\hspace{2.17em}
    如果$\beta$是极限序数, 那么对任意的$\displaystyle
    x<\alpha^\beta=\bigcup_{\gamma<\beta}\alpha^\gamma$, 存在$\gamma<\beta$使得
    \vspace{3pt}
    \[
        x<\alpha^\gamma\quad\implies\quad x^+<\alpha^\gamma+1\leqslant
        \alpha^\gamma2\leqslant\alpha^\gamma\alpha=\alpha^{\gamma^+}
        \quad\implies\quad x^+<\alpha^\beta,
    \]
    所以$\alpha^\beta$是归纳集, 即是极限序数. \hfill $\qedsymbol$
\end{enumerate}





\newpage

\hypertarget{sec:定理1.6}{}
\noindent\textbf{定理1.6 (序数的加法指数)}

任给序数$\alpha,\beta,\gamma$, 有$\alpha^{\beta+\gamma}=\alpha^\beta\alpha^\gamma$.

\noindent{\kaishu 证明.}

给定$\alpha,\beta$, 不妨设$\alpha>1,\beta\neq 0$, 对$\gamma$归纳证明.
任取$\gamma$, 假设对任意的$\delta<\gamma$有
$\alpha^{\beta+\delta}=\alpha^\beta\alpha^\delta$.

\begin{enumerate}
    \item 如果$\gamma=0$, 显然$\alpha^{\beta+0}=\alpha^\beta=\alpha^\beta\alpha^0$.
    
    \vspace{3pt}
    \item 如果$\gamma$是后继序数, 即有$\gamma=\delta^+$, 那么$\delta<\gamma$,
    依假设有
    \[
        \alpha^{\beta+\gamma}=\alpha^{(\beta+\delta)^+}=\alpha^{\beta+\delta}\alpha=
        \alpha^\beta\alpha^\delta\alpha=\alpha^\beta\alpha^\gamma.
    \]

    \vspace{3pt}
    \item 如果$\gamma$是极限序数, 此时$\beta+\gamma$和$\alpha^\gamma$也是极限序数.
    
    \vspace{3pt}\hspace{2.17em}
    一方面, 对任意的$\displaystyle\theta<\beta+\gamma=
    \bigcup_{\delta<\gamma}(\beta+\delta)$, 存在$\delta<\gamma$使得
    $\theta<\beta+\delta$, 所以
    \vspace{3pt}
    \[
        \alpha^\theta<\alpha^{\beta+\delta}=\alpha^\beta\alpha^\delta
        \leqslant\sup_{\rho<\alpha^\gamma}\alpha^\beta\rho
        \quad\implies\quad\sup_{\theta<\beta+\gamma}\alpha^\theta
        \leqslant\sup_{\rho<\alpha^\gamma}\alpha^\beta\rho.
    \]

    \vspace{3pt}\hspace{2.17em}
    另一方面, 对任意的$\displaystyle\rho<\alpha^\gamma=
    \bigcup_{\delta<\gamma}\alpha^\delta$, 存在$\delta<\gamma$使得
    $\rho<\alpha^\delta$, 所以
    \vspace{3pt}
    \[
        \alpha^\beta\rho<\alpha^\beta\alpha^\delta=\alpha^{\beta+\delta}
        \leqslant\sup_{\theta<\beta+\gamma}\alpha^\theta
        \quad\implies\quad\sup_{\rho<\alpha^\gamma}\alpha^\beta\rho
        \leqslant\sup_{\theta<\beta+\gamma}\alpha^\theta.
    \]

    \hspace{2.17em}
    综上即得
    \begin{equation*}
        \alpha^{\beta+\gamma}=\sup_{\theta<\beta+\gamma}\alpha^\theta=
        \sup_{\rho<\alpha^\gamma}\alpha^\beta\rho=\alpha^\beta\alpha^\gamma.
    \tag*{\qedsymbol}\end{equation*}
\end{enumerate}

\vspace{10pt}

\hypertarget{sec:定理1.7}{}
\noindent\textbf{定理1.7 (序数的乘法指数)}

任给序数$\alpha,\beta,\gamma$, 有$\alpha^{\beta\gamma}=(\alpha^\beta)^\gamma$.

\noindent{\kaishu 证明.}

给定$\alpha,\beta$, 不妨设$\alpha>1,\beta\neq 0$, 对$\gamma$归纳证明.
任取$\gamma$, 假设对任意的$\delta<\gamma$有
$\alpha^{\beta\delta}=(\alpha^\beta)^\delta$.

\begin{enumerate}
    \item 如果$\gamma=0$, 显然$\alpha^{\beta0}=1=(\alpha^\beta)^0$.
    
    \vspace{3pt}
    \item 如果$\gamma$是后继序数, 即有$\gamma=\delta^+$, 那么$\delta<\gamma$,
    依假设有
    \[
        \alpha^{\beta\gamma}=\alpha^{\beta\delta+\beta}=
        \alpha^{\beta\delta}\alpha^\beta=(\alpha^\beta)^\delta\alpha^\beta=
        (\alpha^\beta)^\gamma.
    \]

    \vspace{3pt}
    \item 如果$\gamma$是极限序数, 此时$\beta\gamma$也是极限序数.
    
    \vspace{3pt}\hspace{2.17em}
    一方面, 对任意的$\displaystyle\lambda<\beta\gamma=
    \bigcup_{\delta<\gamma}\beta\delta$, 存在$\delta<\gamma$使得
    $\lambda<\beta\delta$, 所以
    \vspace{3pt}
    \[
        \alpha^\lambda<\alpha^{\beta\delta}=(\alpha^\beta)^\delta
        \leqslant\sup_{\delta<\gamma}(\alpha^\beta)^\delta
        \quad\implies\quad\sup_{\lambda<\beta\gamma}\alpha^\lambda
        \leqslant\sup_{\delta<\gamma}(\alpha^\beta)^\delta.
    \]

    \vspace{3pt}\hspace{2.17em}
    另一方面, 对任意的$\delta<\gamma$, 有
    \vspace{3pt}
    \[
        (\alpha^\beta)^\delta=\alpha^{\beta\delta}
        \leqslant\sup_{\lambda<\beta\gamma}\alpha^\lambda
        \quad\implies\quad\sup_{\delta<\gamma}(\alpha^\beta)^\delta
        \leqslant\sup_{\lambda<\beta\gamma}\alpha^\lambda.
    \]

    \hspace{2.17em}
    综上即得
    \begin{equation*}
        \alpha^{\beta\gamma}=\sup_{\lambda<\beta\gamma}\alpha^\lambda=
        \sup_{\delta<\gamma}(\alpha^\beta)^\delta=(\alpha^\beta)^\gamma.
    \tag*{\qedsymbol}\end{equation*}
\end{enumerate}

\vspace{15pt}

最后, 未必有$(\alpha\beta)^\gamma=\alpha^\gamma\beta^\gamma$,
例如$\omega^22^2\neq(\omega2)^2$.

\end{document}